\documentclass[11pt,a4paper]{article}
\usepackage[T1]{fontenc}
\usepackage[utf8]{inputenc}
\usepackage[polish]{babel}
\usepackage{lmodern}
\usepackage{amsmath,amssymb}
\usepackage{geometry}
\usepackage{tikz}
\usetikzlibrary{arrows.meta,positioning,shapes.geometric}
\geometry{margin=2.5cm}

% Style do tworzenia diagramów w TikZiT.
\input{styles.tikzstyles}

\title{Systemy fizyczne jako sieci\\\large Materiał do uzupełnienia diagramami TikZiT}
\author{Programy użytkowe -- ćwiczenie}
\date{}

% Symbol zastępczy. Student zastępuje go poleceniem \input{figures/plik.tikz}.
\newcommand{\missingfigure}[1]{%
  \fbox{\parbox[c][3.2cm][c]{0.83\textwidth}{\centering
  \textbf{MIEJSCE NA RYSUNEK TIKZiT}\\[6pt]#1}}%
}

\begin{document}
\maketitle

\section{Przepływ energii między podukładami}

Różne elementy układu fizycznego mogą przekazywać sobie energię. Rozważmy trzy
części $A$, $B$ i $C$ ułożone w łańcuch. Moc $P_{AB}>0$ płynie od $A$ do $B$,
a moc $P_{BC}>0$ od $B$ do $C$. Bilans energii części środkowej ma postać
\begin{equation}
  \frac{\mathrm{d}E_B}{\mathrm{d}t}=P_{AB}-P_{BC}.
\end{equation}

Diagram powinien pokazywać trzy podpisane węzły i dwa kierunki przepływu,
wraz z etykietami $P_{AB}$ oraz $P_{BC}$. Nie należy odwracać kierunków strzałek.

\begin{figure}[ht]
  \centering
  % TODO: zastąp \missingfigure poleceniem:
  % \input{figures/01_energy_flow.tikz}
  \missingfigure{Trzy węzły $A$, $B$, $C$ i przepływy mocy.}
  \caption{Przepływ energii w układzie trzech elementów.}
  \label{fig:energy}
\end{figure}

\section{Stan równowagi stabilnej}

Rozważmy zmienną $x$ opisującą stan prostego układu dynamicznego:
\begin{equation}
 \dot x=-\gamma\bigl(x-x_{\mathrm{eq}}\bigr),\qquad \gamma>0.
\end{equation}

Gdy $x<x_{\mathrm{eq}}$, pochodna jest dodatnia, więc $x$ rośnie. Gdy
$x>x_{\mathrm{eq}}$, pochodna jest ujemna, więc $x$ maleje. Stan
$x_{\mathrm{eq}}$ jest zatem \emph{stabilny}. W diagramie powinny znaleźć
się trzy pozycje, a strzałki z obu stron powinny prowadzić do stanu środkowego.

\begin{figure}[ht]
  \centering
  % TODO: \input{figures/02_stable_equilibrium.tikz}
  \missingfigure{Stan z lewej i prawej strony prowadzi do równowagi.}
  \caption{Schemat stabilnej równowagi dynamicznej.}
  \label{fig:stable}
\end{figure}

\section{Od pomiaru do wyniku}

Typowy eksperyment wymaga ustalenia kolejnych etapów przetwarzania sygnału.
Zaczynamy od \emph{obiektu badanego}, następnie sygnał trafia do
\emph{detektora}, potem do \emph{analizy danych}, a na końcu otrzymujemy
\emph{wynik wraz z niepewnością}.

Sygnał w idealnym modelu można przedstawić w uproszczonej postaci
\begin{equation}
  y(t)=s(t)+n(t),
\end{equation}

gdzie $s(t)$ jest sygnałem, a $n(t)$ szumem. Diagram powinien pokazywać
cztery etapy skierowane od lewej do prawej. Nie zastępuj rysunku zwykłą listą.

\begin{figure}[ht]
  \centering
  % TODO: \input{figures/03_measurement_chain.tikz}
  \missingfigure{Obiekt $\to$ detektor $\to$ analiza $\to$ wynik.}
  \caption{Łańcuch pomiarowy i opracowanie danych.}
  \label{fig:chain}
\end{figure}

\section*{Podsumowanie}

Uzupełnij trzy miejsca diagramami stworzonymi w TikZiT. Zachowaj spójny wygląd
węzłów i strzałek oraz edytowalne pliki źródłowe \texttt{.tikz}. Sprawdź,
że wszystkie wzory matematyczne i podpisy są czytelne w końcowym PDF-ie.

\end{document}
